\documentclass[a4paper,oneside,14pt]{extreport}

% --- Шрифты и языки (LuaLaTeX) ---
\usepackage{fontspec}
\setmainfont{CMU Serif}
\setsansfont{CMU Sans Serif}
\setmonofont{CMU Typewriter Text}
\usepackage[english,main=russian]{babel}

% --- Поля страницы ---
\usepackage[
    left=30mm,
    right=10mm,
    top=20mm,
    bottom=20mm,
]{geometry}

% --- Переносы ---
\usepackage{microtype}
\sloppy

% --- Межстрочный интервал (полуторный) ---
\usepackage{setspace}
\onehalfspacing

% --- Абзацный отступ ---
\usepackage{indentfirst}
\setlength{\parindent}{12.5mm}

% --- Размеры шрифтов заголовков ---
\makeatletter
\renewcommand\LARGE{\@setfontsize\LARGE{22pt}{26}}
\renewcommand\Large{\@setfontsize\Large{20pt}{24}}
\renewcommand\large{\@setfontsize\large{16pt}{20}}
\makeatother

% --- Настройка заголовков (titlesec) ---
\usepackage{titlesec}
\titleclass{\chapter}{straight}
\titleformat{\chapter}[block]{\hspace{\parindent}\large\bfseries}{\thechapter}{0.5em}{\large\bfseries\raggedright}
\titleformat{name=\chapter,numberless}[block]{\hspace{\parindent}}{}{0pt}{\large\bfseries\centering}
\titlespacing{\chapter}{\parindent}{20pt}{10pt}
\titleformat{\section}[block]{\hspace{\parindent}\large\bfseries}{\thesection}{0.5em}{\large\bfseries\raggedright}
\titleformat{\subsection}[block]{\hspace{\parindent}\large\bfseries}{\thesubsection}{0.5em}{\large\bfseries\raggedright}
\titlespacing{\section}{12.5mm}{10pt}{10pt}
\titlespacing{\subsection}{12.5mm}{10pt}{10pt}

% --- Гиперссылки ---
\usepackage[unicode,hidelinks]{hyperref}

% --- Команда для визуального пропуска (линия для заполнения от руки) ---
\newcommand{\blank}[1][1.5cm]{\underline{\hspace{#1}}}

\usepackage{enumitem}
\makeatletter
\AddEnumerateCounter{\asbuk}{\russian@alph}{щ}
\makeatother
\setlist[enumerate,1]{wide, labelindent=\parindent, labelsep=0.5em, itemsep=2pt, parsep=0pt, topsep=4pt}
\setlist[enumerate,2]{wide, labelindent=\parindent, labelsep=0.5em, itemsep=2pt, parsep=0pt, topsep=4pt}

\begin{document}

% =======================================================================
%  ТИТУЛЬНАЯ СТРАНИЦА ТЗ
% =======================================================================
\begin{titlepage}
    \thispagestyle{empty}

    % --- Верхний блок: СОГЛАСОВАНО / УТВЕРЖДАЮ ---
    \vspace*{0pt}
    \noindent
    \begin{minipage}[t]{0.48\textwidth}
        \raggedright
        \textbf{СОГЛАСОВАНО}\\[6pt]
        Руководитель НИР\\[10pt]
        \blank[3cm]~И.М. Сидякин\\[6pt]
        <<\blank[0.6cm]>>~\blank[2.5cm]~2026~г.
    \end{minipage}%
    \hfill
    \begin{minipage}[t]{0.48\textwidth}
        \raggedright
        \textbf{УТВЕРЖДАЮ}\\[6pt]
        Заведующий кафедрой ИУ3\\[10pt]
        \blank[3cm]~А.Н. Алфимцев\\[6pt]
        <<\blank[0.6cm]>>~\blank[2.5cm]~2026~г.
    \end{minipage}

    % --- Центральный блок: Название ТЗ и тема ---
    \vfill
    \begin{center}
        {\LARGE \bfseries ТЕХНИЧЕСКОЕ ЗАДАНИЕ НА НИР}\\[8mm]
        {\Large \bfseries
        ПРОЕКТИРОВАНИЕ СТРУКТУРИРОВАННОЙ КАБЕЛЬНОЙ\\
        СИСТЕМЫ ОФИСА ЛАБОРАТОРИИ\\
        ТЕСТИРОВАНИЯ УСТРОЙСТВ IoT}
    \end{center}
    \vfill

    % --- Нижний блок: РАЗРАБОТАНО ---
    \noindent
    \textbf{РАЗРАБОТАНО}\\[6pt]
    Студент ИУ3-33Б\\[10pt]
    \blank[3cm]~А.А. Захватов\\[6pt]
    <<\blank[0.6cm]>>~\blank[2.5cm]~2026~г.

    \vspace{20mm}

    % --- Город и год ---
    \begin{center}
        Москва\\
        2026~г.
    \end{center}
\end{titlepage}

\setcounter{page}{2}

% =======================================================================
%  РАЗДЕЛ 1. ОБЩИЕ СВЕДЕНИЯ
% =======================================================================

\chapter{ОБЩИЕ СВЕДЕНИЯ}

\section{Полное наименование системы и ее условное обозначение}
\begin{enumerate}[label=\thesection.\arabic*]
    \item \textbf{Полное наименование:} Структурированная кабельная система (СКС) этажного офиса лаборатории тестирования аппаратуры и программного обеспечения устройств IoT.
    \item \textbf{Условное обозначение:} СКС-IoT-ИУ3-V10.
\end{enumerate}

\section{Шифр темы или основание для проведения работ}
\begin{enumerate}[label=\thesection.\arabic*]
    \item \textbf{Основание для разработки:} Индивидуальное задание к курсовой работе по дисциплине <<Проектирование структурированной кабельной системы>>, 3-й семестр бакалавриата кафедры ИУ-3.
\end{enumerate}

\section{Наименование предприятий (организаций) --- Заказчика и Разработчика}
\begin{enumerate}[label=\thesection.\arabic*]
    \item \textbf{Заказчик:} Кафедра ИУ-3 МГТУ им. Н.Э. Баумана.
    \item \textbf{Разработчик:} Студент группы ИУ3-33Б, Захватов Антон Алексеевич.
\end{enumerate}

\section{Сроки начала и окончания работы}
\begin{enumerate}[label=\thesection.\arabic*]
    \item \textbf{Начало:} Сентябрь 2026 г.
    \item \textbf{Окончание:} Декабрь 2026 г.
\end{enumerate}

% =======================================================================
%  РАЗДЕЛ 2. ЦЕЛИ И НАЗНАЧЕНИЕ СОЗДАНИЯ АВТОМАТИЗИРОВАННОЙ СИСТЕМЫ
% =======================================================================

\chapter{ЦЕЛИ И НАЗНАЧЕНИЕ СОЗДАНИЯ КАБЕЛЬНОЙ СИСТЕМЫ}

\section{Назначение системы}
СКС предназначена для создания универсальной, надежной и масштабируемой кабельной инфраструктуры, обеспечивающей передачу данных, сигналов управления и электропитания по технологии PoE (Power over Ethernet) для обеспечения функционирования:
\begin{enumerate}[label=\asbuk*)]
    \item лабораторных стендов тестирования IoT-устройств и контрольно-измерительного оборудования;
    \item рабочих мест операторов и административного персонала;
    \item систем охранного видеонаблюдения контроля стендов и общих зон;
    \item систем контроля и управления доступом (СКУД);
    \item магистральных каналов связи с серверными здания.
\end{enumerate}

\section{Цели создания системы}
\begin{enumerate}[label=\thesection.\arabic*]
    \item \textbf{Обеспечение требуемой пропускной способности} для передачи трафика от IoT-устройств и систем видеонаблюдения высокого разрешения (в режимах MCU и SFU).
    \item \textbf{Повышение надежности} за счет организации топологии <<иерархическая звезда>> и избыточности магистральных линий.
    \item \textbf{Оптимизация эксплуатации} ИТ-инфраструктуры за счет сведения всех кабельных систем к единому коммутационному центру.
    \item \textbf{Выполнение импортозамещения} путем подбора активного сетевого оборудования из Единого реестра российской радиоэлектронной продукции Минпромторга РФ.
\end{enumerate}

% =======================================================================
%  РАЗДЕЛ 3. ХАРАКТЕРИСТИКА КАБЕЛЬНОЙ СИСТЕМЫ
% =======================================================================

\chapter{ХАРАКТЕРИСТИКА КАБЕЛЬНОЙ СИСТЕМЫ}

\section{Краткие сведения об объекте}
\begin{enumerate}[label=\thesection.\arabic*]
    \item \textbf{Расположение:} Один этаж в многоэтажном здании.
    \item \textbf{Общая площадь:} 613,08~м$^2$ в осях.
    \item \textbf{Количество помещений:} 24 (согласно экспликации архитектурного плана формата А2).
    \item \textbf{Режим работы:} Непрерывный, с повышенной нагрузкой на кабельные тракты в зонах испытаний.
\end{enumerate}

\section{Перечень оборудуемых помещений и плотность рабочих мест}
Распределение информационных портов и оборудования выполняется строго по экспликации:

\begin{enumerate}[label=\thesection.\arabic*]
    \item \textbf{Помещения лабораторий:}
    \begin{enumerate}[label=\thesection.\arabic{enumi}.\arabic*]
        \item \textbf{Помещение №4 (Кабинет 3):} 3~лабораторных стенда. На каждом стенде: 18~портов IoT-устройств, 1~порт камеры контроля (PoE+), 6~резервных портов. Емкость MUTOA стенда = 25~портов.
        \item \textbf{Помещение №11 (Кабинет 1):} 3~лабораторных стенда. На каждом стенде: 8~портов IoT-устройств, 1~порт камеры контроля (PoE+), 6~резервных портов. Емкость MUTOA стенда = 15~портов.
        \item \textit{Дополнительно в №4 и №11:} по 1~камере общего вида (PoE+) и по 1~контроллеру СКУД (PoE+).
    \end{enumerate}

    \item \textbf{Операторская (Помещение №20 <<Опенспейс>>):}
    \begin{enumerate}[label=\thesection.\arabic{enumi}.\arabic*]
        \item \textbf{Пост №1} (обслуживает 1~стенд): 1~моноблок (1~порт PoE++), 1~рабочая станция (1~порт), 1~резервный порт. Итого: 3~розетки.
        \item \textbf{Пост №2} (обслуживает 1~стенд): 1~моноблок (1~порт PoE++), 1~рабочая станция (1~порт), 1~резервный порт. Итого: 3~розетки.
        \item \textbf{Пост №3} (обслуживает 1~стенд): 1~моноблок (1~порт PoE++), 1~рабочая станция (1~порт), 1~резервный порт. Итого: 3~розетки.
        \item \textbf{Пост №4} (обслуживает 1~стенд): 1~моноблок (1~порт PoE++), 1~рабочая станция (1~порт), 1~резервный порт. Итого: 3~розетки.
        \item \textbf{Пост №5} (обслуживает 1~стенд): 1~моноблок (1~порт PoE++), 1~рабочая станция (1~порт), 1~резервный порт. Итого: 3~розетки.
        \item \textbf{Пост №6} (обслуживает 1~стенд): 1~моноблок (1~порт PoE++), 1~рабочая станция (1~порт), 1~резервный порт. Итого: 3~розетки.
        \item \textbf{Общая зона опенспейса:} 2~розетки СКС и 1~штатное место персонала (по 2~розетки).
    \end{enumerate}

    \item \textbf{Административные помещения:}
    \begin{enumerate}[label=\thesection.\arabic{enumi}.\arabic*]
        \item Помещения №24 (Кабинет 2), №7 (Переговорная): оборудуются 2~розетками СКС каждое.
    \end{enumerate}

    \item \textbf{Технологические помещения:}
    \begin{enumerate}[label=\thesection.\arabic{enumi}.\arabic*]
        \item \textbf{Помещение №19:} Серверная этажа (11,70~м$^2$). Размещение главного шкафа.
        \item \textbf{Помещение №16:} Кроссовая СКС (11,70~м$^2$). Размещение этажного распределительного пункта.
    \end{enumerate}
\end{enumerate}

% =======================================================================
%  РАЗДЕЛ 4. ТРЕБОВАНИЯ К АВТОМАТИЗИРОВАННОЙ СИСТЕМЕ
% =======================================================================

\chapter{ТРЕБОВАНИЯ К КАБЕЛЬНОЙ СИСТЕМЕ}

\section{Требования к структуре и архитектуре системы}
\begin{enumerate}[label=\thesection.\arabic*]
    \item СКС должна быть спроектирована по архитектуре \textbf{иерархической звезды} с одним техническим центром (Серверная №19) и одной кроссовой комнатой (Помещение №16).
    \item В кроссовой СКС (№16) должен быть организован Этажный Распределительный Пункт (ЭРП).
    \item Магистральная подсистема должна связывать ЭРП этажа со слаботочным стояком (<<СКС>>) и организовывать вертикальные связи с Главной серверной здания (пом. 903, эт. 9) и Резервной серверной (пом. 1183, эт. 11).
\end{enumerate}

\section{Требования к производительности и компонентам}
\begin{enumerate}[label=\thesection.\arabic*]
    \item Кабельная система горизонтальной подсистемы должна быть выполнена на основе медного четырехпарного кабеля <<витая пара>> категории не ниже \textbf{Cat.6A (класс Ea)} для поддержки скоростей до 10~Гбит/с и обеспечения стабильной работы PoE++/4PPoE (IEEE 802.3bt).
    \item Магистральная подсистема здания (вертикальная) должна быть спроектирована с использованием волоконно-оптического многомодового (OM4) или одномодового (OS2) кабеля для обеспечения отказоустойчивости и емкости каналов.
    \item Коммутационные панели (патч-панели) в шкафах должны иметь плотность не менее 24 портов RJ-45/1U с обязательной цветовой или цифровой маркировкой подсистем.
\end{enumerate}

\section{Требования к надежности и режимам видеонаблюдения}
Система должна обеспечивать бесперебойную передачу трафика от 8~камер (6~направленных, 2~обзорных) с поддержкой кодеков H.265/H.265+ в двух режимах работы медиасерверов:
\begin{enumerate}[label=\asbuk*)]
    \item штатный режим (MCU) --- архитектура с централизованным микшированием потоков;
    \item аварийный режим (SFU) --- архитектура с маршрутизацией сырых потоков напрямую на операторов, требующая проверки пиковых пропускных способностей каналов операторской подсистемы.
\end{enumerate}

\section{Требования к условиям эксплуатации и климатике}
\begin{enumerate}[label=\thesection.\arabic*]
    \item В помещениях Серверной (№19) и Кроссовой (№16) должна быть предусмотрена принудительно-вытяжная вентиляция или кондиционирование для компенсации суммарного тепловыделения активного оборудования.
    \item Параметры микроклимата: температура 18--24~$^\circ$C, относительная влажность 30--55\% без конденсации.
    \item Active-оборудование шкафов должно питаться от выделенных источников бесперебойного питания (ИБП) с топологией On-Line и временем автономной работы при максимальной нагрузке не менее 30 минут.
\end{enumerate}

\section{Требования к кабельным трассам и конструкциям}
\begin{enumerate}[label=\thesection.\arabic*]
    \item Магистральные лотки в коридорах должны быть металлическими (проволочными или перфорированными).
    \item Заполнение лотков и кабельных каналов не должно превышать \textbf{40--50}\% на этапе проектирования с учетом требований по радиусу изгиба и перспективного расширения.
    \item Расчет подвесов и элементов крепления должен проводиться с коэффициентом запаса прочности не менее 3 по отношению к весу заполненного лотка.
\end{enumerate}

\section{Требования к заземлению и электробезопасности}
\begin{enumerate}[label=\thesection.\arabic*]
    \item Все металлические лотки, корпуса шкафов (19"), патч-панели и экраны кабелей должны быть интегрированы в единый контур защитного технологического заземления здания согласно ГОСТ Р 50571.5.54.
    \item Кабели должны отвечать требованиям пожарной безопасности по \textbf{ГОСТ 31565-2012} (исполнение не ниже нг-HF или нг-FRHF).
\end{enumerate}

% =======================================================================
%  РАЗДЕЛ 5. СОСТАВ И СОДЕРЖАНИЕ РАБОТ ПО СОЗДАНИЮ КС
% =======================================================================

\chapter{СОСТАВ И СОДЕРЖАНИЕ РАБОТ ПО СОЗДАНИЮ КАБЕЛЬНОЙ СИСТЕМЫ}
В рамках курсового проектирования (согласно стадиям разработки по ГОСТ Р 15.301-2016) выполняются следующие этапы.
\begin{enumerate}[label=\thechapter.\arabic*]
    \item Анализ исходных данных. Изучение DXF-плана, расчет точного числа портов СКС для каждого помещения.
    \item Расчетно-теоретический этап:
    \begin{enumerate}[label=\asbuk*)]
        \item трассировка и расчет длин кабельных линий (метод контрольных длин);
        \item расчет сечения коробов, лотков, нагрузок на несущие конструкции;
        \item проектирование активных компонентов по реестру Минпромторга (коммутаторы доступа с PoE+/PoE++, агрегаторы, маршрутизаторы);
        \item расчет тепловыделения шкафов, параметров охлаждения и емкости аккумуляторных батарей ИБП.
    \end{enumerate}
    \item Графический этап. Создание схемы размещения рабочих мест на архитектурном плане (А2), составление схемы логической топологии и компоновки шкафов.
\end{enumerate}


% =======================================================================
%  РАЗДЕЛ 6. ПОРЯДОК КОНТРОЛЯ И ПРИЕМКИ КАБЕЛЬНОЙ СИСТЕМЫ
% =======================================================================

\chapter{ПОРЯДОК КОНТРОЛЯ И ПРИЕМКИ КАБЕЛЬНОЙ СИСТЕМЫ}

Контроль результатов учебной разработки осуществляет руководитель НИР.

При контроле проверяются:

\begin{itemize}[label=\asbuk*)]
    \item соответствие проектного решения исходным данным индивидуального варианта;
    \item корректность размещения оборудования и кабельных трасс;
    \item корректность расчета количества портов, патч-панелей и заполнения шкафов;
    \item корректность расчетов трафика в режимах MCU и SFU;
    \item корректность расчетов кабельных линий, лотков, подвесов, охлаждения и ИБП;
    \item обоснованность выбора активного оборудования и PoE-бюджета;
    \item соответствие текстовой и графической документации требованиям задания и применимых нормативных документов.
\end{itemize}

В рамках учебной НИР фактический монтаж и натурные приемочные испытания не выполняются. При реализации проекта кабельные линии подлежат испытаниям и сертификационным измерениям в соответствии с требованиями выбранной кабельной системы и применимых стандартов.

% =======================================================================
%  РАЗДЕЛ 7. ТРЕБОВАНИЯ К СОСТАВУ И СОДЕРЖАНИЮ РАБОТ ПО ПОДГОТОВКЕ
%            ОБЪЕКТА АВТОМАТИЗАЦИИ К ВВОДУ КС В ДЕЙСТВИЕ
% =======================================================================

\chapter{ТРЕБОВАНИЯ К СОСТАВУ И СОДЕРЖАНИЮ РАБОТ ПО ПОДГОТОВКЕ ОБЪЕКТА АВТОМАТИЗАЦИИ К ВВОДУ КАБЕЛЬНОЙ СИСТЕМЫ В ДЕЙСТВИЕ}

Перед вводом кабельной системы в эксплуатацию должны быть выполнены:

\begin{enumerate}[label=\asbuk*)]
    \item подготовка серверной №19 и кроссовой №16 к размещению телекоммуникационных шкафов.
    \item подготовка мест прокладки лотков, коробов и кабельных трасс.
    \item обеспечение электроснабжения, защитного заземления и охлаждения.
    \item монтаж и маркировка информационных розеток, MUTOA, патч-панелей и кабельных линий.
    \item установка и настройка активного оборудования.
    \item проверка кабельных линий и сетевых соединений.
    \item проверка работы PoE+/PoE++ оборудования и резервного питания.
    \item оформление исполнительной документации.
\end{enumerate}

% =======================================================================
%  РАЗДЕЛ 8. ТРЕБОВАНИЯ К ДОКУМЕНТИРОВАНИЮ
% =======================================================================

\chapter{ТРЕБОВАНИЯ К ДОКУМЕНТИРОВАНИЮ}

По результатам НИР должны быть подготовлены:

\begin{enumerate}[label=\asbuk*)]
    \item расчётно-пояснительная записка объёмом от 60 до 100 листов;
    \item техническое задание на кабельную систему;
    \item чертёж схемы сети объекта на листе формата A2;
    \item схема логической топологии сети;
    \item ведомость покупных изделий;
    \item расчёты портов, патч-панелей, шкафов, трафика, кабельных трасс, охлаждения и ИБП.
\end{enumerate}

% =======================================================================
%  РАЗДЕЛ 9. ИСТОЧНИКИ РАЗРАБОТКИ
% =======================================================================

\chapter{ИСТОЧНИКИ РАЗРАБОТКИ}

При разработке используются:

\begin{enumerate}[label=\asbuk*)]
    \item индивидуальное задание на выполнение НИР по теме <<Проектирование структурированной кабельной системы офиса лаборатории тестирования устройств IoT>>;
    \item исходные данные и архитектурный план индивидуального варианта;
    \item DXF-подоснова индивидуального варианта;
    \item ГОСТ Р 70443-2022, устанавливающий требования к техническому заданию на кабельную систему;
    \item ГОСТ 34.602-2020 и ГОСТ Р 15.301-2016, указанные в индивидуальном задании;
    \item ГОСТ 7.32-2017 в части оформления расчетно-пояснительной записки;
    \item ГОСТ Р 53246-2025 и ГОСТ Р 53245-2008 в части проектирования, монтажа и испытаний структурированных кабельных систем;
    \item ГОСТ Р 50571.5.54 в части защитного заземления;
    \item ГОСТ 31565-2012 в части пожарной безопасности кабельных изделий;
    \item Единый реестр российской радиоэлектронной продукции Минпромторга Российской Федерации.
\end{enumerate}


\end{document}